\section{Data and Method}
\label{sec:data-method}

\subsection{Scope and sources}

The study is restricted to the XIV Spanish legislature, oral control questions
addressed to the President of the Government, with answer dates inside the
coverage of ParlaMint 5.0 / ParlaMint-ES (handle 11356/2004; CC BY 4.0) through
23~February 2023. Registered questions are Congreso initiatives of type
\texttt{180/######}. Transcript turns are taken from ParlaMint-ES; registered
wording is taken from Congreso initiative records. Public redistribution of
Congreso full text remains identifier-based until reuse packaging is cleared;
ParlaMint-derived identifiers are releasable under CC BY 4.0 with attribution.

\subsection{Unit of analysis}

The coding unit is the triple formed by (i)~the registered question,
(ii)~the questioner's opening oral turn (Q1), and (iii)~the Prime Minister's
first response (R1). The planned outcome is \emph{reply status} judged on Q1
and anchored by the registered text: explicit reply, partial or intermediate
reply, or absence of reply. Definitions follow a human codebook developed and
calibrated out of pool; the main-study instrument version will be frozen only
after calibration closes (currently draft v0.3 pending Round-2 clarifications).

\subsection{Linkage and attrition}

Candidate PM control exchanges were extracted from ParlaMint-ES (n\,=\,139).
Linkage to Congreso records used sitting date and questioner name against that
day's initiative set, with TEI expediente notes accepted only when
name-confirmed. Link statuses: exact 46, high confidence 54, ambiguous 2,
unlinked 37. Eligible exchanges with complete registered+Q1+R1 text and PM
authorship of R1: n\,=\,100.

A 30-case human linkage QA judged linkage correctness only; no reply-status
coding was performed. Contemporaneous QA records do not name the main
annotators as reviewers. After a Phase~6B contamination audit, those 30 units
remain restored. Three Phase-0 feasibility-probe exchanges were reassessed
under the same rule (linkage feasibility only; no reply-status labels; not
used in out-of-pool development/calibration) and restored. The pre-outcome
main frame is therefore n\,=\,100. No outcome-based exclusions are applied.
Linkage missingness is not assumed random.

\subsection{Formation-vote alignment}

Each questioner is classified by their individual named vote on the XIV second
investiture vote (7~January 2020; DSCD-14-PL-4): supported, opposed, or
abstained. Opposed versus non-opposed (supported $\cup$ abstained) is a
descriptive contextual grouping only, not a confirmatory causal contrast.
Labels such as allies/adversaries are not used analytically.

\subsection{Question-side composition (pre-outcome)}

Registered and Q1 texts in the eligible pool were coded for topic and question
form with deterministic rules, and for an exploratory rule-based
confrontational indicator, without inspecting reply outcomes. The
confrontational indicator is not a validated human measure and is not a
mandatory primary covariate; it documents why causal interpretation of
alignment would be unsafe given a strong imbalance by alignment stratum.
Alignment contrasts in reply explicitness therefore cannot be interpreted as
clean causal effects of political side.

\subsection{Development and calibration}

Instrument development and calibration use out-of-pool prime-ministerial
control exchanges from prior legislatures so that the XIV main frame is not
burned for training. Two expert annotators (Daniel Pinto Pajares; Jose Jaime
Baena Rojas) complete independent coding after a joint clarification step.
Calibration Round~2 is in progress at the time of this writing; main
annotation does not begin until the codebook is frozen.

\subsection{Planned annotation and analysis}

After freeze, both annotators will code the n\,=\,100 main-frame exchanges
independently (same units; distinct deterministic presentation orders).
Agreement is calculated before discussion; original labels are preserved;
adjudicated labels, if used, are stored separately. Descriptive reporting will
include explicit-reply proportions by formation-vote grouping with
questioner-cluster uncertainty and leave-one-questioner-out sensitivity.
Topic strata remain descriptive; the confrontational indicator remains
exploratory/diagnostic only. No LLM labelling is used for the human study.
